\section{Security, QoS and Automation Strategy}\label{sec:sqa}
% Rubric criterion 5 (15 %): micro-segmentation, ACL logic, ZTNA approach, QoS markings for VoIP, scripts explained.

\subsection{Security architecture}\label{sec:security}
The legacy network trusted everything inside its perimeter. The new design applies controls at four
layers, from the access port to the management plane, so that a single failure of policy does not expose
the whole company.

\textbf{1. Access ports.} Every access switch runs DHCP snooping and Dynamic ARP Inspection on the user
VLANs (\cref{lst:l2sec}). Uplinks are trusted; user ports are not. Snooping drops DHCP server messages that
arrive on a user port -- a rogue DHCP server can no longer redirect a VLAN -- and records which IP address
was leased on which port. DAI uses that binding table to drop ARP replies that claim someone else's address,
which defeats ARP-spoofing man-in-the-middle attacks. Each user port also limits the number of MAC
addresses (port security, two sticky addresses, violation mode \texttt{restrict}), shuts down if a switch is
connected (BPDU guard), caps broadcast traffic at 5\,\% of the link (storm control -- the direct answer to
the broadcast storms of the old network), and rate-limits DHCP and ARP to protect the switch CPU.
Option~82 insertion is disabled because the switch is not the relay agent.

\cfgsnippet[lst:l2sec]{l2_access_security.cfg}{Layer-2 security on an access switch}

\textbf{2. Segmentation at the first Layer-3 hop.} Each VLAN's gateway interface has an inbound ACL, so
unwanted traffic is dropped before it is routed anywhere. Because VLAN IDs mean the same thing at every site,
the same policy applies at HQ and at the branches. \Cref{tab:segmentation} shows which zone may open a
connection to which; ``R'' means that only replies pass (TCP \texttt{established} and ICMP echo-reply).

\begin{table}[H]\centering\footnotesize
\caption{Segmentation matrix: who may initiate traffic towards whom}\label{tab:segmentation}
\begin{tabularx}{\textwidth}{@{}>{\raggedright\arraybackslash}p{3.0cm} *{7}{>{\centering\arraybackslash}X}@{}}
\toprule
\textbf{From $\downarrow$ / to $\rightarrow$} & \textbf{Exec / Finance} & \textbf{Other corporate} & \textbf{IoT} & \textbf{DC servers} & \textbf{Mgmt} & \textbf{Guest} & \textbf{Internet} \\
\midrule
Exec / Finance (10)              & --  & yes & yes & yes & no  & no$^{*}$ & yes \\
R\&D, AI lab, Voice (20, 21, 40) & R   & yes & no  & yes & no  & no$^{*}$ & yes \\
Logistics (30)                   & R   & yes & yes & yes & no  & no$^{*}$ & yes \\
Guest (50)                       & no  & no  & no  & no  & no  & --  & yes \\
IoT (60--62)                     & no  & no  & --  & platform only & no & no & no \\
\bottomrule
\end{tabularx}\\[2pt]
\raggedright\footnotesize $^{*}$ Allowed by the source ACL, but guest hosts cannot answer: the guest ACL denies all private destinations.
\end{table}

\Cref{lst:acl} shows the two strictest policies. The guest ACL permits DHCP and HSRP hellos, denies every
RFC~1918 destination and permits the rest -- that is, the Internet only. The IoT ACL is an allow-list:
DHCP, DNS and NTP to the infrastructure server, ping to it for diagnostics, MQTT over TLS (8883) and HTTPS to
the IoT platform, replies to sessions opened by the platform, then an explicit \texttt{deny ip any any log}.
IoT devices are the least patchable systems in the company; with this policy a compromised sensor can reach
its platform and nothing else -- no lateral movement and no Internet. The corporate ACL
(\cref{lst:acl-corp} in the appendix) protects the executive/finance subnets, the management VLANs and the
IoT VLANs of all three sites while allowing everything else. Each IPv4 ACL has an IPv6 twin with the same
logic; the IPv6 versions that end in an explicit deny permit neighbour discovery first, because an explicit
deny removes the implicit ND permits.

\cfgsnippet[lst:acl]{acl_guest_iot.cfg}{Segmentation ACLs for the Guest and IoT VLANs}

\textbf{3. Zone-based firewall at the edge.} The NGFW of the brief is represented by the IOS zone-based
firewall on every edge router (\cref{lst:zbf}). There are three zones: INSIDE (LAN side), OUTSIDE (provider)
and VPN (tunnels). Traffic from INSIDE to OUTSIDE is inspected, so only replies come back and nothing initiated
from the Internet reaches the LAN. Traffic from OUTSIDE to the router itself is limited to what the router
needs: IKE, ESP and GRE, BGP and diagnostic ICMP. The INSIDE--VPN pairs pass traffic statelessly on purpose:
with two edge routers per site a flow can leave through one router and return through the other
(\cref{sec:overlay}), and a stateful rule would drop the returning half because it never saw the session start.

\cfgsnippet[lst:zbf]{zbf.cfg}{Zone-based firewall on the WAN edge}

\textbf{4. Management plane.} Every device uses AAA with a local privileged account whose password is stored
as a hash, accepts SSH version~2 only, and admits management sessions only from the three management VLANs
(\cref{lst:hardening}). Logs go to a central collector on the automation host and carry NTP-synchronised
timestamps; HTTP management is disabled. Automation credentials never appear in a file (\cref{sec:automation}).

\textbf{Limits and the path to Zero Trust.} The ACLs are stateless, so return traffic is admitted by pattern
(\texttt{established}, echo-reply) rather than by session state, and trust is still tied to the VLAN a device
is plugged into. The target state described in the brief is Zero Trust \cite{nist800207}: access is granted
per application, based on the user's identity and the device's posture, and re-evaluated continuously. For TVI
this means 802.1X with RADIUS at the access layer (the identity decides the VLAN and the policy), a stateful
NGFW pair between internal zones, and, for the hybrid workforce, a SASE service: remote users and branch
Internet traffic connect to the nearest SASE point of presence, which provides the secure web gateway, CASB and
ZTNA broker. The local breakout of this design is the first step -- branch Internet traffic would be steered
to the SASE service instead of directly to the provider.

\subsection{Quality of service}\label{sec:qos}
Voice must be prioritised on every WAN link. The design follows the DiffServ model \cite{rfc2474,rfc4594}:
traffic is classified and marked once, at the edge of the network, and every WAN interface queues by that
marking.

\textbf{Trust boundary.} Markings set by end devices are not trusted -- otherwise any PC could promote its
traffic to the voice class. On every LAN-facing interface of an edge router, the input policy
\texttt{PM-MARK-IN} classifies by source and port and rewrites the DSCP value (\cref{tab:qos}); everything
unclassified is reset to best effort. Voice is recognised by its source (the voice VLANs) and the RTP port
range.

\begin{table}[H]\centering\footnotesize
\caption{QoS classes on the WAN}\label{tab:qos}
\begin{tabularx}{\textwidth}{@{}l l l Y@{}}
\toprule
\textbf{Class} & \textbf{Marking} & \textbf{Treatment} & \textbf{Classified by} \\
\midrule
Voice      & EF (46)   & strict priority, 20\,\% & RTP (UDP 16384--32767) from the voice VLANs \\
Video      & AF31 (26) & 25\,\% guaranteed        & UDP 3478--3481 (conferencing media) \\
Signalling & CS3 (24)  & 5\,\% guaranteed         & SIP (5060/5061), SCCP (2000) \\
Network control & CS6 (48) & 5\,\% guaranteed   & routing protocols (marked by the routers) \\
Default    & 0         & fair queueing, remaining bandwidth & everything else \\
\bottomrule
\end{tabularx}
\end{table}

\textbf{Queueing on the WAN.} The output policy on each WAN interface is hierarchical (\cref{lst:qos}). The
parent policy shapes all traffic to the contracted circuit rate, \FactWanRateMbps\,Mbit/s in the lab; the
child policy gives voice a strict-priority queue (low-latency queueing, policed to 20\,\% so that it cannot
starve the other classes \cite{rfc3246}) and guarantees bandwidth to video, signalling and network control.
The shaper is essential: the Ethernet hand-off to the provider is faster than the contracted rate, so without
it the router's queue would never fill, the policy would never act, and the provider would drop packets
without regard to their class. GRE and IPsec copy the DSCP value to the outer header, so the policy still sees
the marking after encryption.

\cfgsnippet[lst:qos]{qos.cfg}{Classification, re-marking and the WAN queueing policy}

\textbf{Sizing.} A G.711 call needs about 87\,kbit/s including Layer-2 overhead, more inside GRE/IPsec. The
priority queue of 20\,\% of 20\,Mbit/s (4\,Mbit/s) therefore carries roughly 40 simultaneous calls. In
production the percentage and the circuit rate must be derived from the expected number of concurrent calls
per site; the structure of the policy stays the same.

\subsection{Network automation}\label{sec:automation}
\textbf{From data to devices.} Automation in this project has two levels. The \emph{build} level generates
everything from the data model: changing a subnet or adding a VLAN in the model and running the generators
updates the startup-configs, the IP plan, the diagrams, the inventory and this report's tables, after which
the static tests check the result. The \emph{operate} level pushes changes to running devices and is the part
the brief specifies in detail.

\textbf{The provisioning script.} \path{Automation/auto_provision.py} uses Netmiko \cite{netmiko} and
implements the three tasks of the brief:

\begin{enumerate}
  \item \textbf{External inventory.} Hosts, management addresses, device types, SSH ports and roles are read
        from \texttt{inventory.yaml} (\cref{lst:inventory}), which is itself generated from the model; the
        script contains no device data.
  \item \textbf{Baseline and VLANs.} The commands come from \texttt{payload.yaml}: AAA, NTP, syslog, SSH
        settings and the banner for every device, plus VLANs 10, 20, 30 and 99 with their names for the
        switches. Additional VLANs can be given on the command line (\texttt{--vlan 110:LAB\_DEMO}). Per-role
        exceptions are data too -- the NTP master, for instance, is not told to synchronise to itself.
  \item \textbf{Backup, verification and logging.} For each device the script captures the output of
        \texttt{show running-config} and of \texttt{show ip interface brief} before and after the change into
        \path{backups/<hostname>_<timestamp>_backup.cfg}, writes a unified diff of the two captures, and
        verifies with \texttt{show vlan brief} that every VLAN of the payload exists. Every step, success or
        failure, is written to \path{logs/automation.log}.
\end{enumerate}

\Cref{lst:process} shows the per-device workflow. Errors are handled per device: an authentication failure,
a connection timeout or any other exception is logged and reported, and the script continues with the next
device; the exit code is non-zero if any device failed, so the script can be used in a pipeline. A
\texttt{--dry-run} prints the commands without connecting, \texttt{--backup-only} changes nothing, and
\texttt{--workers} runs several sessions in parallel.

\pysnippet[lst:process]{py_process_device.py}{Per-device workflow: backup, push, verify, log}

\textbf{No credentials in code.} The script reads the username and passwords from environment variables or
asks for them with \texttt{getpass}, which does not echo (\cref{lst:creds}). The Ansible playbook
\texttt{deploy\_baseline.yml}, the declarative alternative to the script, keeps its credentials in an Ansible
Vault file. \todo{state whether the Ansible playbook was executed in the lab; if not, say so}

\pysnippet[lst:creds]{py_get_credentials.py}{Credentials come from the environment or an interactive prompt}

\textbf{Idempotence and rollback.} Every pushed line describes a desired state (\texttt{ntp server ...},
\texttt{vlan 10} / \texttt{name ...}), so running the script twice changes nothing -- the second run's diff is
empty. The script never removes configuration, and the \emph{before} capture is the rollback.

\textbf{Tests.} \FactTestsAutomation{} unit tests replace the SSH connection with a fake object and check
device selection, that a dry run opens no connection, backup file naming, the handling of timeouts and
authentication failures, VLAN verification, per-role exceptions and the change diff. They run without any
device.

\yamlsnippet[lst:inventory]{inventory_excerpt.yaml}{External inventory (excerpt, generated from the data model)}
